\documentclass[a4paper,11pt,fleqn]{article}%fleqn pour aligner les equations à gauche
\usepackage{../../Styles/Cours}


\newcommand{\chnum}{Chapitre 16}
\newcommand{\chtitle}{Lentilles minces convergentes}
\newcommand{\dtype}{Cours}

\newcommand{\lecon}[1]{
	\colorbox{gray!20}{
		\begin{minipage}{.98\linewidth}
    			\centering \textbf{#1}
		\end{minipage}
	}
}

\begin{document}
\maketitle
\section{L'\oe il }
Un \oe il réel est un récepteur de lumière. Il est constitué de l'\textbf{iris} qui limite la quantité de lumière pénétrant dans l'\oe il par la pupille. L'ensemble des milieux transparents comprenant la \textbf{cornée} et le \textbf{cristallin} forment l'image de l'objet observé sur la rétine. La \textbf{rétine} est située au fond de l'\oe il est là où se forme l'image avant d'être captée et traduite vers les cerveau.\\
L'étude optique de ce système complexe peut être simplifiée en utilisant un modèle : le modèle de l'\oe il réduit.
\begin{table}[h!]
    \centering
    \begin{tabular}{|>{\centering}m{.25\linewidth}|>{\centering}m{.25\linewidth}|>{\centering\arraybackslash}m{.25\linewidth}|}
    \hline  & L'\oe il réel & Modèle de l'\oe il réduit\\
    \hline Limitation de la lumière pénétrant dans l'\oe il & Iris & Diaphragme\\
    \hline Système optique & Ensemble des lumineux transparents (cornée, cristallin ...) & Lentille mince convergente\\
    \hline Lieu de formation de l'image & Rétine & Écran situé à une distance constante $d$ de la lentille\\
    \hline 
    \end{tabular}
\end{table}

\section{Lentilles minces convergentes}
    \subsection{Modèle du rayon lumineux}
    Un \textbf{rayon lumineux} modélise le trajet suivi par la lumière pour aller d'un point à un autre. Il est représenté par un segment de droite fléché, la flèche indiquant le sens de propagation. C'est un \textbf{modèle :} il n'a pas de réalité physique. Une source émet toujours un faisceau lumineux constitué d'une infinité de rayons.

    \subsection{Modèle de la lentille mince convergente}
    Une lentille est dite minci si son épaisseur $e$ au centre est très inférieure aux rayons de courbure $R_1$ et $R_2$ de ses faces. Elle est convergente si le bord est plus mince que le centre.\\
    Le modèle d'une lentille mince convergente néglige l'épaisseur de la lentille. La lentille mince convergente est représentée par un segment fléché.
    \begin{itemize}
        \item l'\textbf{axe optique} est l'axe de symétrie de la lentille.
        \item le centre de la lentille est appelé \textbf{centre optique} $O$.
    \end{itemize}
    \lecon{Tout rayon lumineux passant par le centre optique d'une lentille mince convergente n'est pas dévié.}
    \begin{itemize}
        \item Les foyers $F$ et $F'$ d'une lentille convergente sont deux points particuliers de l'axe optique.
    \end{itemize}
    $F$ et $F'$ sont symétriques par rapport au centre optique $O$. La distance focale $f'$ d'une lentille mince convergente est la distance qui sépare le centre optique $O$ du foyer image $F'$ : $$f'=OF'$$
    \lecon{\begin{itemize}
        \item Tout rayon lumineux incident parallèle à l'axe optique émerge d'une lentille mince convergents en passant par le foyer image F'.
       \item Tout rayon lumineux incident passant par le foyer objet F émerge d'une lentille minci convergente parallèlement à l'axe optique.
    \end{itemize}}

\section{Image réelle d'un objet réel formée par une lentille mince convergente}
    \subsection{Point objet et point image}
    \begin{itemize}
        \item Un \textbf{objet réel} est modélisé par un ensemble de points objets. Un point objet est le point de croisement des rayons lumineux qui arrivent sur la lentille.
        \item Un \textbf{point image} est le point de croisement des rayons lumineux qui émergent de la lentille. À chaque point objet correspond un point image et un seul. 
    \end{itemize}

    \subsection{Construction de l'image réelle d'un objet réel}
        \subsubsection{Objet réel $AB$ situé à une distance finie de la lentille}
        Le point $B'$ image du point objet $B$ est construit à l'aide d'au moins deux des trois rayons lumineux dont le chemin est connu.\\
        Le point $A'$ image du point objet $a$ est tel que $A'B'$ est perpendiculaire à l'axe optique.

        \subsubsection{Objet réel $AB$ situé à l'infini}
        Lorsqu'un objet $AB$ est à l'infini, le faisceau issu de chaque point de cet objet et arrivant sur la lentille est un faisceau parallèle.\\
        Le point image $A'$ est alors confondu avec le foyer image $F'$. Le point image $B'$ est dans le plan orthogonal à l'axe optique passant par $F'$, appelé plan focal image.\\ Pour construire le point image $B'$, le rayon lumineux passant par le centre optique suffit.
    \subsection{Grandissement}
    Pour comparer la taille et le sens de l'image $A'B'$ à ceux de l'objet $AB$, on détermine le \textbf{grandissement} $\mathbf{\gamma}$.\\
    
    \lecon{\begin{itemize}
        \item Si l'image est dans le même sens que l'objet : $\mathbf{\gamma = \dfrac{A'B'}{AB}}$
        \item Si l'image et l'objet sont dans des sens opposés : $\mathbf{\gamma = -\dfrac{A'B'}{AB}}$
    \end{itemize}}
\end{document}